\section{Complejidad y contexto}
El problema en cuestión esta contenido en la clase de problemas NP-completo. Se argumenta su pertenecía a la clase NP y su NP completitud.

\subsection{Pertenencia a \textit{NP}}

Un problema pertenece a la clase \textit{NP} si, dadas las soluciones, es posible verificar si estas soluciones son correctas en tiempo polinominal. El problema se puede ejemplificar suponiendo que tenemos como solución un conjunto de \textit{K} subconjuntos, que son los que particionan el conjunto de los vértices \textit{V}. Para verificar que una solución dada cumple con la propiedad de solución del problema, la solución debe cumplir las siguientes condiciones:

\begin{itemize}
    \item Verificar que cada vértice pertenece a exactamente un subconjunto.
    \item Para cada subconjunto, comprobar que todos los pares de vértices están conectados entre sí, es decir, que cada subconjunto produce una clique.
\end{itemize}

La verificación de la primera condición se lleva a cabo en tiempo polinominal, porque consiste en comprobar si las aristas entre los pares de vértices dentro de cada subconjunto.

Por todo lo anterior se asume que el problema de \textit{Clique Cover} pertenece a la clase \textit{NP}.

\subsection{\textit{NP}-completitud}

Para demostrar que el problema es \textit{NP-completo} hay que comprobar que es \textit{NP-duro}. Esto se hace llevando a cabo una reducción polinomial desde un problema que se sabe que es \textit{NP-completo}. 

En esta ocasión se considerará el problema de \textit{Graph Coloring}. Este problema consiste en asignar colores a los vértices de un grafo de manera que no existan dos vértices adyacentes que estén coloreados del mismo color, y que ademas estén representados con el menor numero de colores posibles. 

La demostración se realizara considerando la relación existente entre un grafo $G$ y su grafo complemento $\overline{G}$, el cual esta definido sobre el mismo conjunto de vértices, donde existe una arista entre dos vértices si y solo si dicha arista no pertenece a $E$.

En este contexto, se observa que: 

\begin{itemize}
    \item Un conjunto de vértices forma un clique $G$ si y solo si se forma un conjunto independiente en $\overline{G}$.
    \item Cubrir los vértices de $G$ con $K$ cliques es equivalente a colorear el grafo $\overline{G}$ con sus $K$ colores.
\end{itemize}

Dado que el problema \textit{Graph Coloring} es \textit{NP-completo} y existe una transformacion polinomial y el problema de \textit{Clique Cover}, se concluye que \textit{Clique Cover} en \textit{NP-duro}.

Como aparte pertenece a \textit{NP} es posible llegar a la conclusión de que el problema es \textit{NP-completo}.


\subsection{Relación con otros problemas}

El problema abordado tiene una relación estrecha con otros problemas de la teoría de grafos en estos se incluyen \textit{Vertex Cover} y \textit{Graph Coloring}. Estos problemas forman la parte fundamental del conjunto de problemas de \textit{NP-completos} y comparten estructuras t transformaciones similares.  

\begin{table}[h]
\centering
\caption{Relación entre Clique Cover, Vertex Cover y Graph Coloring}
\begin{tabular}{|l|p{4cm}|p{3.5cm}|p{5cm}|}
\hline
\textbf{Problema} & \textbf{Definición} & \textbf{Objetivo} & \textbf{Relación con otros} \\ \hline

Clique Cover 
& Particionar los vértices en cliques (subgrafos completos) 
& Minimizar el número de cliques 
& Equivalente a Graph Coloring en el grafo complemento \\ \hline

Vertex Cover 
& Conjunto de vértices que cubre todas las aristas 
& Minimizar el número de vértices 
& Complementario de Independent Set \\ \hline

Graph Coloring 
& Asignar colores a vértices sin que adyacentes compartan color 
& Minimizar el número de colores 
& Equivalente a Clique Cover en el complemento \\ \hline

\end{tabular}
\end{table}


